\scanpage{26}
\uline{Foncteurs distingués} : \inserted{Sommes disjointes finies de} quotient d'un foncteur
\deleted{\unsure{spécial par.}} $h_I$ par un groupe d'autom, i.e.
foncteurs isomorphes à des \inserted{sommes finies de} foncteurs de la
forme
\[
E\longmapsto\operatorname{Mon}(I,E)/H
\]
où $H\subset\mathfrak S_I$. Ce sont aussi les sommes
finies de composantes connexes de foncteurs
de la forme $H(I,E)/H=\operatorname{Hom}_{\mathrm{Ens}}(I,E)/H$.
Ces foncteurs ont les \unsure{$=^{s}$|mêmes} propriétés de
permanence que celles vues pour les
foncteurs spéciaux.

\margin{Tout foncteur \uline{distingué} est
isom. à un foncteur
de la forme
\[
E\longmapsto\coprod_{n\geq0}\operatorname{Mon}(I_n,E)\mathbin{\wedge}_{\mathfrak S_n}X_n
\]
où $\forall n$, $X_n$ est un $\mathfrak S_n$-ens. fini (\unsure{à g.}).}

Un hom \inserted{$u:F\longrightarrow G$} de foncteurs $F,G:\mathcal E\longrightarrow(\mathrm{Ens})$
est dit une équivalence s'il satisfait la
condition suivante

($\unsure{\sigma|\alpha}$) $\exists n\in\mathbf N$ tel que pour \deleted{deux} tout
ensemble fini $E$ \deleted{$\unsure{F}$} de cardinal $\geq n$, \deleted{et les \unsure{appl.}}
\deleted{\unsure{inj.} $E\hookrightarrow F$ \unsure{on a} $\unsure{F(E)\simeq G(E)}$,}
on a $u(E):F(E)\simeq G(E)$. On a

a) Si $F\mathrel{\overset{u}{\text{---}}}G\mathrel{\overset{v}{\text{---}}}H$ et si deux parmi les
trois $u,v,vu$ sont des équivalences, aussi le $3^\circ$.

b) Les \deleted{\unsure{morphismes}} équivalences sont stables par
chgt de base
\mspage{25}
